% 原书分章PDF第4—6页：21.1.2—21.1.4，所有干预定义/目标与直接效应。
\subsection{干预类型}

感兴趣的干预通过将条件分布 $\bm G$ 替换为用户提供的新干预 $\bm G^*$ 来定义；假定 $\bm G^*$ 具有定义良好的密度 $\bm g^*$。也就是说，$\bm G^*$ 是一个多元条件分布，规定了如何在给定 $\bm W$ 时生成每个干预后的暴露，记为 $A_i^*$。注意，对 $\bm A$ 的静态或动态干预，对应于 $\bm g^*$ 的退化选择；而 $\bm g^*$ 的非退化选择通常称为随机干预（例如 Dawid 和 Didelez，2010）。假定 $\bm A$ 和 $\bm A^*$ 属于同一个共同空间 $\mathcal A$，不对 $\bm G^*$ 作进一步限制。另定义 $A_i^{*s}:=a_i^s(\bm A^*)$，其中 $A_i^{*s}$ 表示由汇总量映射作用于干预暴露向量 $\bm A^*$ 所蕴含的随机变量，$i=1,\ldots,N$。最后，将 $P_{\bm Q,\bm G}$ 中的 $\bm G$ 替换为用户提供的新分布 $\bm G^*$，定义干预后分布 $P_{\bm Q,\bm G^*}$。用 $\bm O^*=(W_i,A_i^*,Y_i^*)_{i=1}^N$ 表示在 $P_{\bm Q,\bm G^*}$ 下生成的随机变量，其似然为
\begin{equation}
 p_{\bm Q,\bm G^*}^N(\bm O^*)
 =\left[\prod_{i=1}^N q_Y(Y_i^*\mid A_i^{*s},W_i^s)\right]
  \bm g^*(\bm A^*\mid\bm W)\bm q_{\bm W}(\bm W).
 \tag{21.4}
\end{equation}
后一个分布 $P_{\bm Q,\bm G^*}$ 称为干预 $\bm G^*$ 下 $\bm O$ 的干预后分布的 g 计算公式（Robins，1987）。

\subsection{目标参数：结局期望的样本平均}

将目标统计量 $\psi_0$ 定义为式（21.4）中干预后分布的函数。具体而言，
\begin{equation}
 \psi_0=\Psi(P_0^N)
 =\frac{1}{N}\sum_{i=1}^N E_{\bm q_0,\bm g^*}[Y_i^*],
 \tag{21.5}
\end{equation}
即相依单位 $i=1,\ldots,N$ 的 $N$ 个结局的样本平均的期望，期望相对于干预后分布 $P_{\bm Q,\bm G^*}$ 计算。我们将 $\Psi(P_0^N)$ 视为从统计模型 $\mathcal M$ 到 $\RR$ 的映射。$\psi_0$ 是在给定观测网络结构 $\bm F$ 的条件下定义的，也以 $N$ 为索引。另定义 $\bar Q(A_i^s,W_i^s)=\int_y yq_Y(y\mid A_i^s,W_i^s)\,d\mu(y)$，表示在对所有 $i$ 共同的分布 $Q_Y$ 下计算的条件均值；$\bar Q_0$ 则表示在 $P_0^N$ 处取值的 $\bar Q$。降维假设意味着 $E_{P_0^N}[Y_i\mid\bm A,\bm W]=\bar Q_0(A_i^s,W_i^s)$。由于目标参数 $\psi_0$ 仅通过 $\bar Q_0$ 与 $\bm Q_{\bm W,0}$ 依赖于 $P_0^N$，略微滥用记号，我们将交替使用 $\Psi(P_0^N)$ 与 $\Psi(\bar Q_0,\bm Q_{\bm W,0})$。因此，参数 $\psi_0$ 以 $N$、$\bm F$ 和 $\bm G^*$ 为索引，可以写为
\begin{equation}
 \psi_0=\frac{1}{N}\sum_{i=1}^N\int_{\bm a,\bm w}
 \bar Q_0(a_i^s(\bm a,\bm w),w_i^s(\bm w))
 \bm g^*(\bm a\mid\bm w)\bm q_{\bm W,0}(\bm w)\,d\mu(\bm a,\bm w),
 \tag{21.6}
\end{equation}
其中积分相对于某个控制测度 $\mu(\bm a,\bm w)$ 进行。

对于许多实践者而言，把目标量定义为两个随机干预的对比可能很有价值。例如，可以将 $\Psi^{\bm G_1^*}(P_0^N)$ 和 $\Psi^{\bm G_2^*}(P_0^N)$ 分别定义为在随机干预 $\bm G_1^*$ 与 $\bm G_2^*$ 下计算的上述目标参数。于是，可以定义目标对比 $\Psi^{\bm G_1^*,\bm G_2^*}(P_0^N)=\Psi^{\bm G_1^*}(P_0^N)-\Psi^{\bm G_2^*}(P_0^N)$。$N$ 个相互连接单位的平均处理效应是 $\Psi^{\bm G_1^*,\bm G_2^*}(P_0^N)$ 的一个特例，对应于退化的 $\bm G_1^*$、$\bm G_2^*$：对所有 $\bm w\in\mathcal W$，由 $\bm g_1^*(\bm 1^N\mid\bm w)=1$ 和 $\bm g_2^*(\bm 0^N\mid\bm w)=1$ 定义。也就是说，$\bm G_1^*$ 将所有研究对象的处理都设为1，$\bm G_2^*$ 则将其全部设为0。下文只关注为某个特定 $\bm G^*$ 定义的统计参数 $\psi_0$ 的估计，并指出所有结果都可自然推广到这类对比。

要把 $\psi_0$ 解释为一个因果量，即在干预 $\bm G^*$ 下，$N$ 个单位组成的网络中反事实结局期望的样本平均，除A1—A4、B1和C1以外，还需额外的因果假设（见 van der Laan，2014a）。这些因果假设不对上述概率分布 $P_0^N$ 施加进一步限制，统计模型 $\mathcal M$ 保持不变。由于 $\mathcal M$ 包含真实数据分布 $P_0^N$，$\psi_0$ 总具有纯粹的统计解释，即数据分布 $P_0^N$ 的特征 $\Psi(P_0^N)$。因果假设在当前估计问题中不起作用：即使不相信任何不可检验的因果假设，统计参数 $\psi_0$ 仍然是在控制全部已测量混杂因素后，一个感兴趣的效应度量。最后，只需选择一个以观测基线协变量 $\bm W$ 为条件的目标参数 $\psi_0$，就可以完全避开假设A1。关于由此得到的估计问题，van der Laan（2014a）有深入描述。

\subsection{个体间干扰下的样本平均均值直接效应}

下面说明，如何方便地推广目标参数的定义，以一组单位 $i$ 特定的随机干预 $\{g_{F_i^*}^*:i=1,\ldots,N\}$ 为索引，定义结局期望的样本平均。此时，每个 $g_{F_i^*}^*$ 都表示对 $i$ 的暴露 $A_i$ 以及 $i$ 的朋友的暴露 $(A_j:j\in F_i)$ 的干预。\footnote{注意，前面已将 $F_i^*$ 定义为 $F_i\cup\{i\}$。}例如，考虑在 $N$ 名相互连接的个体组成的网络中，存在个体间干扰时，估计二元暴露直接效应的问题。在存在个体间干扰时，直接效应将单位 $i$ 的暴露对其自身结局的影响，与 $i$ 的朋友的暴露可能对 $i$ 的结局产生的影响分离开来。为刻画直接效应，可以通过把单位自身的暴露 $A_i$ 设为0或1，同时将所有 $j\in F_i$ 的 $A_j$ 保持在观测值 $(a_j:j\in F_i)$，来定义各个 $g_{F_i^*}^*$。另一种选择是，按照朋友暴露的观测条件分布 $G_0(A_j\mid W_j^s)$，对 $(A_j:j\in F_i)$ 进行随机干预。

每个结局 $Y_i^{g_{F_i^*}^*}$ 均从它自身、特定于 $i$ 的干预后分布生成；这一分布用 $g_{F_i^*}^*$ 替换 $i$ 及 $j\in F_i$ 的观测处理分配。于是，结局 $Y_i^{g_{F_i^*}^*}$ 的期望的样本平均为
\begin{equation}
 \psi_0=\frac{1}{N}\sum_{i=1}^N
 E_{\bm q_0,g_{F_i^*}^*}\left[Y_i^{g_{F_i^*}^*}\right].
 \tag{21.7}
\end{equation}
参数 $\psi_0$ 也可以表达为
\begin{equation}
 \psi_0=\frac{1}{N}\sum_{i=1}^N\int_{\bm a,\bm w}
 \bar Q_0(a_i^s(\bm a,\bm w),w_i^s(\bm w))
 g_{F_i^*}^*(\bm a\mid\bm w)\bm q_{\bm W,0}(\bm w)\,d\mu(\bm a,\bm w).
 \tag{21.8}
\end{equation}
因此，个体间干扰下的直接效应可以表示为两个目标参数的对比。第一个参数以一组干预为索引，这些干预将每个 $A_i$ 设为1，但保持每组暴露 $(A_j:j\in F_i)$ 不变。第二个参数则以一组将每个 $A_i$ 设为0、同时保持每组 $(A_j:j\in F_i)$ 不变的干预为索引。注意，可以调整后续各节给出的框架来估计这些参数；更多细节见 Sofrygin 和 van der Laan（2017）第8节。
